\chapter{Consigna 1}
\label{chap:diseno}

\section{Diseño del Filtro}

Se adopta la topología de la Figura~\ref{fig:circuito}: la resistencia $R$ y el capacitor $C$ en serie con
la entrada, y el inductor $L$ en paralelo sobre la salida. La elección se fundamenta en que, con esta
disposición, la resistencia interna del generador $R_g$ queda en serie con $R$ y se incorpora al diseño
como una resistencia total $R_t = R_g + R$, de modo que una resistencia de generador comparativamente baja
no modifica apreciablemente el comportamiento del filtro.

\begin{figure}[!ht]
  \centering
  \begin{subfigure}[t]{0.48\textwidth}
    \centering
    \begin{tikzpicture}[american]
  \draw (0,0) to[short, o-] (0.4,0)
        to[R, l=$R$] (2.2,0)
        to[C, l=$C$] (4,0)
        to[short, -o] (5.5,0);
  \draw (4,0) to[L, l=$L$, *-*] (4,-2.5);
  \draw (0,-2.5) to[short, o-o] (5.5,-2.5);
  \draw (4,-2.5) node[ground]{};
  \draw (0,0) to[open, v=$v_i(t)$] (0,-2.5);
  \draw (5.5,0) to[open, v=$v_o(t)$] (5.5,-2.5);
\end{tikzpicture}

    \caption{Circuito en el dominio del tiempo.}
    \label{crkt:pasa-alto}
  \end{subfigure}
  \hfill
  \begin{subfigure}[t]{0.48\textwidth}
    \centering
    \begin{tikzpicture}[american]
  \draw (0,0) to[short, o-] (0.4,0)
        to[generic, l=$R$, i=$I(s)$] (2.2,0)
        to[generic, l=\tiny$\frac{1}{sC}$] (4,0)
        to[short, -o] (5.5,0);
  \draw (4,0) to[generic, l=$sL$, *-*] (4,-2.5);
  \draw (0,-2.5) to[short, o-o] (5.5,-2.5);
  \draw (4,-2.5) node[ground]{};
  \draw (0,0) to[open, v=$V_i(s)$] (0,-2.5);
  \draw (5.5,0) to[open, v=$V_o(s)$] (5.5,-2.5);
\end{tikzpicture}

    \caption{Modelo de impedancias en el dominio de Laplace.}
    \label{crkt:laplace}
  \end{subfigure}
  \caption{Filtro pasa alto pasivo RLC de segundo orden.}
  \label{fig:circuito}
\end{figure}

\section{Función de transferencia}

Analizando el circuito de la Figura~\ref{crkt:laplace}, la corriente de malla es

\begin{equation}
  I(s) = \frac{V_i(s)}{R + \dfrac{1}{sC} + sL} = \frac{sC}{s^2 LC + sRC + 1}\, V_i(s),
  \label{eq:corriente}
\end{equation}

y la tensión de salida es la caída sobre el inductor, $V_o(s) = sL\, I(s)$. Reemplazando la
ecuación~\eqref{eq:corriente} y dividiendo numerador y denominador por $LC$ se obtiene la función de
transferencia en términos de los componentes:

\begin{equation}
  H(s) = \frac{V_o(s)}{V_i(s)} = \frac{s^2 LC}{s^2 LC + sRC + 1}
       = \frac{s^2}{s^2 + s\,\dfrac{R}{L} + \dfrac{1}{LC}}.
  \label{eq:tf-componentes}
\end{equation}

La ecuación~\eqref{eq:tf-componentes} presenta un cero doble en el origen y dos polos dados por las raíces
del polinomio característico, lo que confirma el comportamiento pasa alto de segundo orden: para $s \to 0$
la ganancia tiende a cero y para $s \to \infty$ tiende a la unidad.

\section{Forma canónica y ecuaciones de diseño}

La forma canónica de un pasa alto de segundo orden es

\begin{equation}
  H(s) = \frac{s^2}{s^2 + \dfrac{\omega_0}{Q}\, s + \omega_0^2},
  \label{eq:tf-canonica}
\end{equation}

donde $\omega_0$ es la pulsación natural y $Q$ el factor de calidad. Comparando término a término las
ecuaciones~\eqref{eq:tf-componentes} y~\eqref{eq:tf-canonica} resulta el sistema

\begin{gather*}
  \omega_0^2 = \frac{1}{LC}, \qquad \frac{\omega_0}{Q} = \frac{R}{L},
\end{gather*}

de donde, fijando $R$ como parámetro libre, se despejan los valores de los componentes:

\begin{equation}
  L = \frac{R\,Q}{\omega_0}, \qquad C = \frac{1}{\omega_0\, R\, Q}.
  \label{eq:diseno-LC}
\end{equation}

\subsection{Relación entre la frecuencia de corte y la pulsación natural}

La pulsación natural $\omega_0$ coincide con la pulsación de corte $\omega_c$ (atenuación de
\qty{-3}{\decibel}) únicamente cuando $Q = 1/\sqrt{2}$ (respuesta de Butterworth). Para cualquier otro valor
de $Q$ es necesario introducir una constante de corrección $K$ tal que $\omega_0 = K\,\omega_c$. Su
expresión se obtiene imponiendo la condición de \qty{-3}{\decibel} sobre el módulo de la
ecuación~\eqref{eq:tf-canonica}. Definiendo $x = \omega/\omega_0$,

\begin{gather*}
  |H(jx)|^2 = \frac{x^4}{\left(1 - x^2\right)^2 + \dfrac{x^2}{Q^2}} = \frac{1}{2}
  \;\Longrightarrow\;
  x^4 + \left(2 - \frac{1}{Q^2}\right) x^2 - 1 = 0,
\end{gather*}

cuya única solución positiva, llamando $a = 1 - \dfrac{1}{2Q^2}$, es $x_c^2 = -a + \sqrt{a^2 + 1}$. Como
$K = \omega_0/\omega_c = 1/x_c$, racionalizando se obtiene

\begin{equation}
  K = \sqrt{\left(1 - \frac{1}{2Q^2}\right) + \sqrt{\left(1 - \frac{1}{2Q^2}\right)^2 + 1}}.
  \label{eq:factor-K}
\end{equation}

Reemplazando $\omega_0 = 2\pi f_c K$ en la ecuación~\eqref{eq:diseno-LC}, las ecuaciones de diseño
quedan expresadas directamente en función de los datos del problema:

\begin{equation}
  L = \frac{R\,Q}{2\pi f_c\, K}, \qquad C = \frac{1}{2\pi f_c\, K\, Q\, R}.
  \label{eq:diseno-final}
\end{equation}

\subsection{Factor de amortiguamiento}

Expresando el denominador de la ecuación~\eqref{eq:tf-canonica} en la forma estándar de los sistemas de
segundo orden, $s^2 + 2\xi\omega_0 s + \omega_0^2$, el factor de amortiguamiento resulta

\begin{equation}
  \xi = \frac{1}{2Q},
  \label{eq:xi}
\end{equation}

y su valor determina el tipo de amortiguamiento: sobreamortiguado si $\xi > 1$ (raíces reales y distintas),
críticamente amortiguado si $\xi = 1$ (raíz real doble) y subamortiguado si $\xi < 1$ (raíces complejas
conjugadas).

\section{Elección de la resistencia}

Se elige $R = \qty{100}{\ohm}$ como compromiso entre dos criterios: un valor suficientemente alto para que
la resistencia interna del generador (típicamente \qty{50}{\ohm} o menor) no altere de manera significativa
el factor de calidad del filtro, y suficientemente bajo para que, según la ecuación~\eqref{eq:diseno-final},
no resulten valores de inductancia excesivamente elevados.

\section{Valores de los componentes}

Aplicando las ecuaciones~\eqref{eq:factor-K}, \eqref{eq:diseno-final} y~\eqref{eq:xi} con
$f_c = \qty{400}{\hertz}$ y $R = \qty{100}{\ohm}$ se obtienen los valores de la
Tabla~\ref{tab:componentes}. A modo de ejemplo, para $Q = \num{0.25}$:

\begin{gather*}
  K = \sqrt{\left(1 - \frac{1}{2 \cdot \num{0.25}^2}\right)
      + \sqrt{\left(1 - \frac{1}{2 \cdot \num{0.25}^2}\right)^2 + 1}} = \num{0.2666}, \\[4pt]
  L = \frac{\qty{100}{\ohm} \cdot \num{0.25}}{2\pi \cdot \qty{400}{\hertz} \cdot \num{0.2666}}
    = \qty{37.31}{\milli\henry}, \qquad
  C = \frac{1}{2\pi \cdot \qty{400}{\hertz} \cdot \num{0.2666} \cdot \num{0.25} \cdot \qty{100}{\ohm}}
    = \qty{59.70}{\micro\farad}.
\end{gather*}

\begin{gather*}
  K = \sqrt{\left(1 - \frac{1}{2 \cdot \num{0.75}^2}\right)
      + \sqrt{\left(1 - \frac{1}{2 \cdot \num{0.75}^2}\right)^2 + 1}} = \num{1.057}, \\[4pt]
  L = \frac{\qty{100}{\ohm} \cdot \num{0.75}}{2\pi \cdot \qty{400}{\hertz} \cdot \num{1.057}}
    = \qty{28.23}{\milli\henry}, \qquad
  C = \frac{1}{2\pi \cdot \qty{400}{\hertz} \cdot \num{1.057} \cdot \num{0.75} \cdot \qty{100}{\ohm}}
    = \qty{5.02}{\micro\farad}.
\end{gather*}

\begin{table}[H]
  \centering
  \begin{tabular}{|c|c|c|c|c|c|c|}
    \hline
    $Q$ & $K$ & $\omega_0$ [\unit{\radian\per\second}] & $f_0$ [\unit{\hertz}] & $L$ & $C$ & $\xi$ \\
    \hline
    \num{0.25} & \num{0.2666} & \num{670.0}  & \num{106.6} & \qty{37.31}{\milli\henry}
               & \qty{59.70}{\micro\farad} & \num{2.000} \\
    \num{0.75} & \num{1.0570} & \num{2656.5} & \num{422.8} & \qty{28.23}{\milli\henry}
               & \qty{5.02}{\micro\farad}  & \num{0.667} \\
    \hline
  \end{tabular}
  \caption{Constante de corrección, pulsación natural, componentes y factor de amortiguamiento de cada diseño.}
  \label{tab:componentes}
\end{table}

